% ============================================================
% NITheCS Candidate's Statement — Prof. Muaaz Bhamjee
% Candidate's Research Vision and Prospective Contributions
% ============================================================
\documentclass[11pt,a4paper]{article}

\usepackage[a4paper, top=2.2cm, bottom=2.2cm, left=2.5cm, right=2.5cm]{geometry}
\usepackage{hyperref}
\usepackage{enumitem}
\usepackage{titlesec}
\usepackage{parskip}
\usepackage[T1]{fontenc}
\usepackage{xcolor}
\usepackage{booktabs}
\usepackage{array}
\usepackage{tabularx}

% ── Colour palette ────────────────────────────────────────
\definecolor{nithecs}{RGB}{0, 70, 127}
\definecolor{rulecolour}{RGB}{0, 70, 127}

% ── Section formatting ────────────────────────────────────
\titleformat{\section}
  {\large\bfseries\color{nithecs}}
  {}{0em}{}[\color{rulecolour}\titlerule]

\titleformat{\subsection}
  {\normalsize\bfseries\color{nithecs}}
  {}{0em}{}

\titlespacing{\section}{0pt}{14pt}{6pt}
\titlespacing{\subsection}{0pt}{10pt}{4pt}

% ── Hyperlinks ────────────────────────────────────────────
\hypersetup{
  colorlinks=true,
  urlcolor=nithecs,
  linkcolor=nithecs,
  pdfauthor={Muaaz Bhamjee},
  pdftitle={NITheCS Candidate Statement — Prof. Muaaz Bhamjee},
}

% ── List spacing ─────────────────────────────────────────
\setlist[itemize]{noitemsep, topsep=4pt, leftmargin=1.5em}
\setlist[enumerate]{noitemsep, topsep=4pt, leftmargin=1.5em}

% ── No paragraph indent ──────────────────────────────────
\setlength{\parindent}{0pt}
\setlength{\parskip}{6pt}

% =============================================================
\begin{document}

% ── Header ───────────────────────────────────────────────────
{\Large\bfseries Prof.\ Muaaz Bhamjee}\\[4pt]
{\normalsize
  Associate Professor, Department of Mechanical \& Aeronautical Engineering\\
  Faculty of Engineering, Built Environment and Information Technology\\ 
  University of Pretoria \\
  \href{mailto:muaaz.bhamjee@up.ac.za}{muaaz.bhamjee@up.ac.za} | \qquad
  \href{Office Contact: +27 12 420 5366}{+27 12 420 5366} | \qquad
  \href{Mobile +27 76 474 4434}{+27 76 474 4434} \\
  \href{URL}{https://muaazbhamjee.github.io} | \qquad
  \href{URL}{https://github.com/muaazbhamjee} \qquad
}

\vspace{4pt}
\noindent\color{rulecolour}\rule{\linewidth}{1.5pt}\color{black}

\vspace{6pt}
\begin{center}
  {\large\bfseries Candidate's Statement}\\[3pt]
  {\normalsize\itshape Research Vision and Prospective Contributions to the National Institute\\
  for Theoretical and Computational Sciences (NITheCS)}
\end{center}
\vspace{2pt}
\noindent\color{rulecolour}\rule{\linewidth}{0.5pt}\color{black}
\vspace{6pt}

% ── 1. INTRODUCTION ──────────────────────────────────────────
\section{Introduction}

I am Prof.\ Muaaz Bhamjee, Associate Professor in the Department of Mechanical and Aeronautical Engineering at the University of Pretoria. 
My research programme spans computational fluid dynamics and granular mechanics, machine learning and AI applied to fluid dynamics, geospatial and climate systems, 
experimental high-energy physics through the ATLAS Collaboration at CERN, and emerging quantum computing and quantum sensing programmes that I am establishing at UP. 
I hold a BIng and MIng (Mechanical Engineering), a BSc Honours (Applied Mathematics), and a DIng (Computational Fluid and Granular Dynamics), all from the University of Johannesburg, 
where I served as a Lectuer and Senior Lecturer before transitioning to IBM Research Africa as a Staff Research Scientist and subsequently to my current position at UP.

The unifying thread across all of these programmes is \textbf{physics-informed computational simulation of transport phenomena}, the development and validation of numerical models 
that resolve how physical systems transport energy, matter, and information, applied wherever the physics demands it: from granular particles in a 
hydrocyclone to acoustic waves in a diseased lung to Higgs Bosons decaying in the ATLAS detector.

In 2026 I established CERG-FLUX Lab (Fluids, Learning, and Uncertainty in compleX systems) as a named research subgroup within the Clean Energy Research Group at the University of Pretoria, 
providing a structured home for my research programme across CFD, Lattice Boltzmann Methods, Scientific Machine Learning, ATLAS/CERN particle physics, and quantum technologies. 
The group currently comprises 3 MEng students, 6 PhD students, and postdoctoral fellows in pipeline, with an active digital presence and open-source research infrastructure

I apply to become a NITheCS Affiliated Fellow with the conviction that the institute's 
transdisciplinary model uniting theoretical, computational, and applied sciences across all 26 South African public 
universities is precisely the environment in which my research vision can achieve its broadest societal impact. 
The University of Pretoria provides the infrastructure, postgraduate pipeline, and industry partnerships to ground this vision nationally, while NITheCS 
provides the network to scale it.

% ── 2. CURRENT RESEARCH PORTFOLIO ────────────────────────────
\section{Current Research Portfolio and Alignment with NITheCS Themes}

\subsection{2.1\quad Computational Fluid Dynamics, Granular Mechanics and Multiphase Flow}

My primary research programme, which forms the core of my academic group at UP, develops and applies the \textbf{Lattice Boltzmann Method (LBM)}, Euler--Euler and Euler-Largange CFD frameworks to resolve multiphase and granular flow in industrial process equipment. This includes:

\begin{itemize}
  \item Hydrocyclone dynamics under surging conditions, published in \textit{Mathematical and Computational Applications} (2022) and the \textit{R\&D Journal of the South African Institution of Mechanical Engineering} (2021), with industrial partner Multotec Pty.\ Ltd.
  \item Heated gas-solid fluidised bed modelling using Eulerian frameworks (2021), and solar thermal CFD investigations published in \textit{Solar Energy} (2020) and \textit{Energy and Buildings} (2013, IF 5.879).
  \item An LBM--DEM solver currently under active development targeting dense granular flow regimes inaccessible to continuum approaches.
  \item Biomedical CFD: a UV germicidal irradiation droplet infectiousness model published in \textit{Computers \& Fluids} (2024); ventriculoperitoneal shunt dynamics; pulmonary acoustics work co-supervised with students and published in \textit{Biomedical Signal Processing and Control} (2026).
\end{itemize}

This programme aligns with NITheCS's \textbf{Complex Systems and Applied Mathematics} themes and connects to engineering industry through Multotec and the CHPC's industrial HPC initiative.

\subsection{2.2\quad Machine Learning, AI, and Earth Systems / Climate Modelling}

Building on work initiated at IBM Research Africa, I continue to develop AI-based geospatial intelligence tools with direct climate relevance:

\begin{itemize}
  \item Co-development and open-source release of \textit{Granite Geospatial Land Surface Temperature}, an Earth Observation Foundation Model fine-tuned from IBM's \textit{Prithvi} Vision Transformer, hosted on HuggingFace.
  \item Development of \textit{Auxiliary Feature Injection} for ViT fine-tuning on high-resolution geospatial data, presented at IEEE IGARSS 2025 (Brisbane).
  \item Machine-learning detection and characterisation of Urban Heat Islands in African cities, presented at IEEE IGARSS 2024 (Athens) and the Deep Learning Indaba 2024.
  \item Reservoir computing for predicting chaotic dynamical systems, presented at the 69th SAIP Annual Conference 2025, extending the AI thread into fundamental physics data pipelines.
  \item A granted South African patent (ZA2023/09448) for dynamically forecasting high-resolution air temperature using multi-source sensor fusion.
\end{itemize}

These outputs align with NITheCS's \textbf{Earth Systems and Climate Change Modelling} and \textbf{Machine Learning and AI} flagships, and address the DSTI 2022--32 Decadal Plan's Africa-centred climate adaptation priority.

\subsection{2.3\quad High-Energy Physics: ATLAS at CERN}

I am a contributing author on the \textbf{ATLAS Collaboration} at CERN, with authorship across more than 200 publications (Google Scholar h-index: 53; Scopus h-index: 25):

\begin{itemize}
  \item Co-authorship of the landmark \textit{Nature} (2024) paper reporting the \textbf{first observation of quantum entanglement with top quarks}, a defining result of Run~2 and one of the most cited experimental physics papers of 2024.
  \item Co-authorship of a comprehensive \textit{Physics Reports} (2025) review synthesising the full ATLAS Run~2 Higgs characterisation programme.
  \item Environmental monitoring for the Inner Tracker (ITk) upgrade in preparation for the High-Luminosity LHC (HL-LHC), and supervision of two candidates through their ATLAS Authorship Qualification Tasks (AQTs).
\end{itemize}

This work connects to NITheCS's \textbf{Theoretical Physics}, \textbf{Quantum Technologies}, and \textbf{Gravity, Astrophysics and Cosmology} themes.

\subsection{2.4\quad Quantum Computing and Quantum Sensing}

New research programmes in quantum computing and quantum sensing are being established at the University of Pretoria. This is enabled by the University of Pretoria South African Qunatum Technology Initiative (SA QuTI) node, University of Pretoria Quatum Science and Technology (UPQust) hub and  emerging industry partnerships. Initial research directions include quantum algorithm development for simulation of fluid-mechanical systems, and quantum sensor architectures for mineral separation.

% ── 3. RESEARCH VISION ───────────────────────────────────────
\section{Research Vision as a NITheCS Fellow}

My vision as a NITheCS Affiliated Fellow is to build a \textbf{Pan-African Computational and Quantum Research Laboratory} at UP---a programme that:

\begin{enumerate}
  \item \textbf{Develops LBM-DEM frameworks for modeeling process equipment} Develop LBM-DEM framework for modelling process equipment. Embed CFD/LBM physics as constraints in machine learning architectures to accelerate such models whilst preserving the key mechanics/physics. 

  \item \textbf{Bridges HEP data science and engineering AI.} The ATLAS HL-LHC programme will generate data volumes that demand advanced ML. The pipelines developed for detector monitoring and event reconstruction---including reservoir computing and transformer architectures---transfer directly to industrial process optimisation and geospatial analytics. NITheCS provides the network to formalise and scale this transfer.

  \item \textbf{Establishes quantum computing and sensing as engineering research tools.} Using UP's quantum programmes as a foundation, I will develop quantum simulation frameworks for granular and multiphase flow---a class of problems that remains computationally intractable for classical HPC at the densest flow regimes---and quantum sensing arrays for mineral separation. NITheCS's connections with SA QuTI will be central to this.

  \item \textbf{Scales HPC-enabled research across the NITheCS network.} Drawing on my CHPC experience---annual conference presenter since 2015, Invited speaker in 2017, 2020 and 2025; I will develop shared HPC workflows for LBM, ML, Quantum Computing, and geospatial analytics accessible across the NITheCS node network.

  \item \textbf{Advances quantitative climate-health research.} My CFD droplet infectiousness and Lung accoustic research form the nucleus of a health research thread engaging NITheCS's Afrocentric Health Options flagship and collaborators at the Perinatal HIV Research Unit (PHRU) at Wits.
\end{enumerate}

% ── 4. NATIONAL AND INTERNATIONAL COLLABORATIONS ─────────────
\section{Existing Collaborations Supporting the NITheCS Programme}

\begin{tabularx}{\linewidth}{lX}
\toprule
\textbf{Partner} & \textbf{Nature of Collaboration} \\
\midrule
CERN / ATLAS & Co-author on 200+ papers; AQT supervisor; ITk environmental monitoring \\
IBM Research (Global) & Co-developer of open-source Granite Geospatial EOFM; patent co-inventor \\
CHPC (SA) & Annual presenter; HPC school co-organiser; LBM and ML workflow enabler \\
Multotec Pty.\ Ltd. & Industry research collaboration; hydrocyclone and granular process CFD \\
ESRF (European Synchrotron) & BEATS beamline detector mechanics \\
PHRU / Wits & COVID-19 droplet modelling; respiratory health CFD \\
University of Johannesburg & Ongoing postgraduate co-supervision; SA-CERN node activities \\
\bottomrule
\end{tabularx}

% ── 5. SUPERVISION AND TRAINING ──────────────────────────────
\section{Supervision, Training and Capacity Development}

\begin{itemize}
  \item Supervised 9 MEng students to completion (3 graduated Cum Laude); currently supervising/co-supervising 3 MEng/MSc and 6 PhD candidates at UP and UJ.
  \item Supervised 2 ATLAS Authorship Qualification Tasks (one completed March 2022), connecting South African students to CERN's global research environment.
  \item Mentored postdoctoral fellows and research interns at IBM Research Africa.
\end{itemize}

% ── 6. COMMUNITY SERVICE AND GOVERNANCE ──────────────────────
\section{Community Engagement and Scientific Leadership}

\begin{itemize}
  \item \textbf{SAAM Vice-President}: Vice-President of the Executive Committee of the South African Association for Theoretical and Applied Mechanics (SAAM ExCo).
  \item \textbf{National IUTAM Vice-President}: Vice-President of the South African National Committee of the International Union of Theoretical and Applied Mechanics.
\end{itemize}

% ── 7. CONCLUSION ────────────────────────────────────────────
\section{Conclusion}

NITheCS represents South Africa's boldest investment in transdisciplinary theoretical and computational science. Its mandate, to unite disparate disciplines and make them more than the sum of their parts, resonates directly with my research trajectory: from LBM granular mechanics to quantum entanglement at CERN, from CFD respiratory models to AI-powered climate intelligence, from industrial process optimisation to quantum sensing for mineral processing.

From my position at the University of Pretoria---a research-intensive institution with strong engineering, physics, and computational science capacity---I bring a \textbf{productive, internationally connected, and impact-oriented research programme} that is simultaneously rooted in Africa's development agenda and competitive at the global frontier. As a NITheCS Fellow, I am committed to contributing to the colloquium and school programme, co-supervising students across the network, contributing to HPC enablement, and building the interdisciplinary bridges---between engineering mechanics, particle physics, AI, and quantum technologies---that NITheCS was created to make possible.

\vspace{20pt}
\noindent\textit{Submitted by:}

\vspace{8pt}
\noindent\textbf{Prof.\ Muaaz Bhamjee}\\
PrEng, DIng (Mech); MIng (Mech); BSc Honours (Applied Mathematics); BIng (Mech) \\
Associate Professor: Clean Energy Research Group, Department of Mechanical and Aeronautical Engineering, Faculty of Engineering, the Built Environment and Imnformation Technology \\
University of Pretoria\\
University of Pretoria ATLAS CERN Institutional Representative/Team Leader \\
Tel:  +27 (0)12 420 5366 \\
Room 6-63,  Engineering 3 Building \\
University of Pretoria, Private Bag X20, Hatfield 0028, South Africa \\
Email: \href{mailto:muaaz.bhamjee@up.ac.za}{muaaz.bhamjee@up.ac.za} \\

----------- \\
Community Engagement \\
Vice-President of the South African Association for Theoretical and Applied Mechanics (SAAM): \url{https://saam.africa/exco/} \\
Vice-President of the South African National Committee of the International Union of Theoretical and Applied Mechanics (IUTAM) \\
----------- \\
\end{document}
